% !TEX program = xelatex
% Pointing as Morphology -- the neural/ companion paper.
% Build:  xelatex main && bibtex main && xelatex main && xelatex main
% Requires XeLaTeX + a Syriac-capable font (Noto Sans Syriac). If that font is
% absent, comment out the \newfontfamily\syriacfont line; the \syr macro then
% falls back to the main font (Syriac may not render, transliteration still does).

\documentclass[11pt]{article}

\usepackage[margin=1in]{geometry}
\usepackage{amsmath,amssymb}
\usepackage{booktabs}
\usepackage{array}
\usepackage{multirow}
\usepackage{microtype}
\usepackage[round]{natbib}
\usepackage[hidelinks]{hyperref}
\usepackage{xcolor}

% --- Fonts / Syriac support (XeLaTeX) ------------------------------------ %
\usepackage{fontspec}
% Right-to-left Syriac font. If "Noto Sans Syriac" is not installed, comment the
% next line out; the macro below then falls back to the main font, and the
% accompanying transliteration always renders regardless.
\newfontfamily\syriacfont[Script=Syriac]{Noto Sans Syriac}
\providecommand{\syriacfont}{}            % no-op fallback if the line above is removed
\providecommand{\RL}[1]{#1}               % overwritten by the bidi package below

% \syr{script}{transliteration}{gloss}: native Syriac + translit + English,
% so readers who cannot read the script always have an anchor.
\newcommand{\syrscript}[1]{{\syriacfont\RL{#1}}}
\newcommand{\syr}[3]{\syrscript{#1}~(\textit{#2}, `#3')}
\newcommand{\syrt}[2]{\syrscript{#1}~(\textit{#2})}

\newcommand{\auc}{\textsc{auc}}
\newcommand{\cal}[1]{\texttt{#1}}         % CAL-ASCII transliteration of skeletons
\newcommand{\todo}[1]{\textcolor{red}{[#1]}}

% The bidi package must be loaded LAST (after hyperref and other packages).
\usepackage{bidi}

\title{Pointing as Morphology:\\
A Self-Supervised Vocaliser and Cross-Script Transfer\\
for Classical Syriac}

\author{Anonymous\\ \texttt{TODO: affiliation} \and Anonymous\\ \texttt{TODO: affiliation}}
\date{\today}

\begin{document}
\maketitle

\begin{abstract}
Classical Syriac has no large pretrained language model, and from-scratch
pretraining is ruled out by scale: the openly digitised corpus is on the order of
a few million tokens, one to two orders of magnitude short of what monolingual
pretraining needs. We therefore ask how to build useful neural representations for
Syriac by injecting information from sources other than more Syriac text, and we
report what works and what does not on the same authorship cohort and metric as a
prior count-based study. Our central result is a vocaliser that treats the
restoration of Syriac pointing as a morphological task: in a defective abjad the
unwritten vowels are the \emph{pattern} morpheme of a root-and-pattern system, so
predicting them is recovering morphology, not orthographic decoration. Trained on
the SEDRA lexicon, the model restores held-out pointing at $0.825$ per slot. We
then show it generalises to a different register: although the Digital Syriac
Corpus was assumed to be consonantal, $56\%$ of its tokens in fact carry vowel
points, which we turn into a $134{,}000$-form classical gold set by deriving the
mark-to-vowel mapping from data; on it the lexicon-trained vocaliser reaches
$0.638\pm0.028$ vowel accuracy on in-vocabulary forms and $0.594\pm0.019$ on forms
whose consonantal skeleton it never saw. For document-level authorship we compare
three transfer routes for a zero-resource abjad: a tokenizer-free byte model
(CANINE) reaches \auc{} $0.870$ off the shelf; transliterating Syriac into Hebrew
script and encoding it with a Hebrew model reaches $0.888$; and a larger
multilingual model whose tokenizer represents Syriac as character fallback reaches
only $0.798$ -- so script coverage, not language count, is what matters. A
supervised verification head lifts either encoder to $\approx 0.96\pm0.03$, on par
with the count-based baselines rather than beating them and seed-sensitive on the
small author pool: the neural representations match, but do not beat, those
baselines at document-level authorship, so the value of the transfer comparison is
the mechanism it exposes, not a new best.
Continued pretraining improves the intrinsic language-model objective but not
document-level authorship; an explicit root/pattern factorisation does not beat a
flat encoder when the consonantal skeleton is already visible. Code, the
evaluation harness, and the gold-set construction are released; the SEDRA-derived
data is regenerated locally under its licence rather than redistributed.
\end{abstract}

% ===================================================================== %
\section{Introduction}
\label{sec:intro}

Classical Syriac---a dialect of Aramaic and a principal literary language of late
antique Christianity---has a large, partly anonymous or contested corpus and very
little modern language technology. There is no large pretrained Syriac language
model, and the reason is not architectural but quantitative: a from-scratch
monolingual Transformer expects on the order of $10^{8}$--$10^{9}$ tokens, whereas
the Digital Syriac Corpus (DSC) holds about $2.18$M and even aggregating every
openly digitised Syriac text reaches only the low tens of millions. Pretraining
from scratch is therefore the wrong tool, and so is the tempting shortcut of
generating ``more Syriac'' from a model fit on Syriac, which adds no information
and risks degenerate feedback \citep{shumailov2024collapse}.

If new information cannot come from more raw Syriac, it must come from elsewhere.
We organise the design around three such sources. \textbf{Structure}: the
root-and-pattern grammar can be supplied as explicit supervision. \textbf{Transfer}:
knowledge can be carried over from byte- or character-level models and from related
Semitic languages. \textbf{Real text}: the genuinely new tokens that do exist
(here, a vocalisation signal that the corpus turns out to contain) can be put to
work. This paper develops one method in each category and evaluates them against
the cohort and metric of a prior count-based study of the same corpus, so that the
neural numbers are directly comparable.

The work is deliberately a study of \emph{methods and findings} rather than of new
architectures: the components---a tokenizer-free encoder \citep{clark2022canine},
low-rank adaptation \citep{hu2022lora}, a Hebrew model \citep{seker2022alephbert},
a supervised-contrastive head \citep{khosla2020supcon}, a sequence labeller---are
all established. The contribution is what they reveal about Syriac and, more
generally, about modelling a zero-resource Semitic abjad.

\paragraph{Contributions.}
\begin{enumerate}
  \item \textbf{Pointing as morphology.} We frame vocalisation restoration as
  recovery of the root-and-pattern \emph{pattern} morpheme and build a vocaliser
  for Classical Syriac (\S\ref{sec:vocaliser}). To our knowledge it is the first
  neural Syriac vocaliser; diacritisation itself is established for Arabic and
  Hebrew \citep{zalmout2019arabic,gershuni2022hebrew}, so the novelty is the
  language and the morphological framing.
  \item \textbf{A classical vocalisation gold set.} We show the DSC is $56\%$
  vocalised, derive the mapping from its Unicode vowel points to a phonemic label
  set directly from data, and release the construction of a $134{,}000$-form
  classical gold set on which cross-register transfer can be measured
  (\S\ref{sec:vocaliser}).
  \item \textbf{Transfer routes for an abjad.} We compare a tokenizer-free byte
  encoder, cross-script transfer through Hebrew, and a large multilingual model on
  identical authorship evaluation, and explain the ordering by tokenizer coverage
  (\S\ref{sec:transfer}).
  \item \textbf{What does and does not help.} A supervised verification head gives
  the best authorship separation but is seed-sensitive on a small author pool;
  continued pretraining helps the intrinsic objective but not authorship; and an
  explicit root/pattern factorisation does not beat a flat encoder
  (\S\ref{sec:transfer},~\S\ref{sec:other}). We report these as results, with the
  conditions under which each holds.
\end{enumerate}

% ===================================================================== %
\section{Related Work}
\label{sec:related}

\paragraph{Tokenizer-free and multilingual encoders.}
Subword vocabularies built for high-resource languages cover an abjad like Syriac
poorly. Tokenizer-free models avoid the problem by reading bytes or codepoints:
CANINE encodes Unicode codepoints directly \citep{clark2022canine} and ByT5
operates on bytes \citep{xue2022byt5}, both extending character-aware modelling
\citep{kim2016character} to the Transformer \citep{vaswani2017attention}.
Massively multilingual subword models such as XLM-R \citep{conneau2020xlmr} and
Glot500 \citep{imani2023glot500} cover hundreds of languages, but nominal coverage
of a script does not guarantee useful subwords for it, as we show.

\paragraph{Cross-lingual transfer and adaptation.}
When target data is scarce, knowledge is transferred from related languages or
adapted with few parameters. Continued in-domain pretraining \citep{gururangan2020dapt}
and parameter-efficient methods---adapters \citep{houlsby2019adapters} and
low-rank adaptation \citep{hu2022lora}---are the standard tools at our data scale.
Syriac and Hebrew are both $22$-letter Aramaic-derived abjads, which motivates
transliterating Syriac into Hebrew script to reuse a Hebrew model
\citep{seker2022alephbert}; a Syriac morphological parser has likewise borrowed
Hebrew data to offset scarcity \citep{naaijer2023transformer}.

\paragraph{Diacritisation and vocalisation.}
Restoring vowels to a consonantal script is well studied for Arabic
\citep{zalmout2019arabic} and Hebrew \citep{gershuni2022hebrew}, typically as a
sequence-labelling or reading-aid task. We reuse that machinery but reframe the
objective: for a root-and-pattern language the pointing \emph{is} the pattern
morpheme, so vocalisation is morphological self-supervision rather than a
downstream convenience.

\paragraph{Authorship and restoration.}
Burrows's Delta \citep{burrows2002delta} and the broader attribution literature
\citep{stamatatos2009survey} provide the stylometric baselines; a
supervised-contrastive head \citep{khosla2020supcon} provides a learned metric. We
assess uncertainty with the author-cluster bootstrap \citep{efron1986bootstrap}.
For application, neural restoration of damaged ancient text is established for
Greek and Latin \citep{assael2022ithaca}; we include a Syriac character-level
restorer as a demonstrator. The authorship questions are long-standing in
philology; the corpus attributed to Ephrem, for instance, has been reassessed on
text-critical grounds \citep{hartung2018authorship}. For the language see
\citet{brock2006introduction,butts2019syriac}; the data is the DSC
\citep{digitalsyriaccorpus}, ETCBC Syriac \citep{etcbcsyriac}, and the SEDRA
lexicon \citep{kiraz1994sedra}; tooling is \textsc{PyTorch}
\citep{paszke2019pytorch} and \textsc{Transformers} \citep{wolf2020transformers}.

% ===================================================================== %
\section{Data and Resources}
\label{sec:data}

\paragraph{Corpora.}
The DSC (the \texttt{srophe/syriac-corpus} TEI/XML release) provides $632$
authored texts and $2.18$M tokens of classical prose and verse; it is our running
text and, below, our classical vocalisation gold. The ETCBC Syriac New Testament
and Peshitta \citep{etcbcsyriac} provide additional consonantal running text with
morphological annotation. The SEDRA~3 lexicon \citep{kiraz1994sedra} provides
$29{,}699$ word forms with explicit vocalisation, root, and morphology; it is the
supervision for the vocaliser. SEDRA is distributed for academic use without
redistribution of altered versions, so we ship the code that regenerates the
derived data from a user-provided SEDRA source rather than the data itself, and we
cite Kiraz as required.

\paragraph{The CAL skeleton alphabet and transliteration.}
SEDRA encodes forms in a CAL-style ASCII alphabet that is \emph{non-phonetic}: the
letter is a fixed slot in the Aramaic alphabet order, so for instance \cal{K} is
\emph{Heth}, \cal{C} is \emph{Kaph}, \cal{;} is \emph{Yudh}, \cal{/} is
\emph{Sadhe}, and \cal{X} is \emph{Qaph}. We add a deterministic Syriac$\leftrightarrow$CAL
map (and a Syriac$\to$Hebrew map for the transfer experiment); all $29{,}699$
SEDRA skeletons round-trip exactly through it. A vocalised SEDRA form interleaves
lowercase vowels and diacritics with the consonantal skeleton, e.g.\
\cal{AaB,oHaOH\_;} has skeleton \cal{ABHOH;} (\syrt{ܐܒܗܘܗܝ}{abāhaw(hy)}).

\paragraph{The corpus is vocalised.}
The DSC was assumed to be consonantal, and the prior count-based study strips
combining marks. In fact $55.9\%$ of DSC tokens---across $600$ of $632$
files---carry Syriac vowel points (U+0730--U+074A: \emph{pthaha}, \emph{zqapha},
\emph{rbasa}, \emph{hbasa}, \emph{esasa}, the East-Syriac dotted \emph{zlama}s,
and \emph{qu\v{s}\v{s}aya}/\emph{rukkaka}). This is a genuine classical
vocalisation signal, distinct in register from SEDRA's New-Testament lexicon, and
\S\ref{sec:vocaliser} turns it into an evaluation gold set.

% ===================================================================== %
\section{Pointing as Morphology: A Syriac Vocaliser}
\label{sec:vocaliser}

\paragraph{Objective.}
A Syriac word is a consonantal root interleaved with a vocalic pattern; the script
writes the consonants and usually omits the pattern. Restoring the pointing is
therefore recovering the pattern morpheme. We frame it as sequence labelling: the
input is the consonant skeleton, and for each consonant slot the model predicts
the vowel/diacritic string that follows it (one of $30$ classes including
``bare''). Every position is supervised, so the model cannot collapse to the prior
the way a sparsely-masked objective can on small data. The model is a two-layer
bidirectional LSTM ($603$K parameters), the standard diacritiser architecture.

\paragraph{Held-out lexicon accuracy.}
On a $90/10$ type-level split of SEDRA, the vocaliser restores $0.825$ of slots
correctly (per-consonant majority baseline $0.507$) and reproduces $0.397$ of full
words exactly (baseline $0.003$); see the upper block of Table~\ref{tab:vocaliser}.
These are held-out \emph{New-Testament-vocabulary} forms, which raises the obvious
question of register.

\paragraph{Cross-register evaluation on classical text.}
We answer it with the DSC vocalisation signal. Because the corpus marks vowels with
Unicode points and SEDRA marks them with CAL letters, we need a mapping between the
two; rather than assert one, we \emph{derive} it from data. Aligning the $541{,}416$
DSC tokens whose consonantal skeleton has a unique SEDRA vocalisation, we tally,
per consonant slot, which Unicode mark co-occurs with which CAL vowel. The vowels
map cleanly---\emph{pthaha}$\to$\cal{a}, \emph{zqapha}$\to$\cal{o},
\emph{rbasa}$\to$\cal{e}, \emph{hbasa}$\to$\cal{i}, \emph{esasa}$\to$\cal{u} at
$0.90$--$0.98$ purity---while reading and grammatical marks fall below threshold and
are treated as unvocalised (Table~\ref{tab:vowelmap}). Applying the derived map to
all vocalised DSC tokens yields a gold set of $133{,}960$ unique
(skeleton, pattern) forms, of which $47{,}973$ have a skeleton in SEDRA's
vocabulary and $85{,}987$ do not.

Classical scribes point selectively, so the fair metric is vowel accuracy on the
slots that are actually pointed. The lexicon-trained vocaliser, applied unchanged
to classical text, reaches $0.638\pm0.028$ on in-vocabulary skeletons and
$0.594\pm0.019$ on out-of-vocabulary skeletons---forms whose consonantal skeleton
never appeared in training (mean $\pm$ SD over five seeds; per-consonant baselines
$0.40$ and $0.45$; lower block of Table~\ref{tab:vocaliser}). The small gap between
the in- and out-of-vocabulary splits indicates the model has learned the
root-and-pattern regularity rather than memorised lexical entries, which is the
property that makes vocalisation a useful morphological objective.

\begin{table}[t]
  \centering
  % Generated by neural/results.py -- do not edit by hand.
\begin{tabular}{lcc}
\toprule
 & vowel/pointing acc. & baseline \\
\midrule
\multicolumn{3}{l}{\emph{Held-out SEDRA (NT lexicon)}} \\
\quad per-position pointing & 0.825 & 0.507 \\
\quad full-word exact match & 0.397 & 0.003 \\
\midrule
\multicolumn{3}{l}{\emph{Classical DSC, 5 seeds}} \\
\quad in SEDRA vocabulary & 0.638$\pm$0.028 & 0.403 \\
\quad out of SEDRA vocab. & 0.594$\pm$0.019 & 0.447 \\
\bottomrule
\end{tabular}

  \caption{The vocaliser on held-out SEDRA lexicon forms (upper) and on classical
  DSC text (lower). Cross-register figures are vowel accuracy on pointed slots,
  mean $\pm$ SD over five seeds; the out-of-vocabulary split contains consonantal
  skeletons unseen in training.}
  \label{tab:vocaliser}
\end{table}

\begin{table}[t]
  \centering
  % Generated by neural/results.py -- do not edit by hand.
\begin{tabular}{llrr}
\toprule
Syriac mark (U+0730--074A) & $\to$ CAL vowel & purity & count \\
\midrule
Zqapha Dotted$^\ast$ & o & 0.85 & 122,748 \\
Pthaha Dotted$^\ast$ & a & 0.90 & 122,652 \\
Pthaha Above$^\ast$ & a & 0.97 & 102,975 \\
Zqapha Above$^\ast$ & o & 0.97 & 85,556 \\
Rbasa Above$^\ast$ & e & 0.96 & 72,855 \\
Dotted Zlama Horizontal$^\ast$ & e & 0.90 & 47,265 \\
Esasa Above$^\ast$ & u & 0.98 & 41,127 \\
Dotted Zlama Angular$^\ast$ & e & 0.78 & 29,055 \\
Hbasa Above$^\ast$ & i & 0.96 & 28,988 \\
Feminine Dot & u & 0.38 & 8,074 \\
\bottomrule
\end{tabular}

  \caption{The Unicode-point$\to$CAL-vowel mapping \emph{derived from data} by
  aligning DSC tokens to uniquely-vocalised SEDRA skeletons. Marks accepted as
  vowels are starred; reading and grammatical marks (not shown) fall below the
  majority threshold and are treated as unvocalised.}
  \label{tab:vowelmap}
\end{table}

% ===================================================================== %
\section{Transfer for a Zero-Resource Abjad}
\label{sec:transfer}

We now turn to document-level authorship. The evaluation is the prior study's:
centred-cosine separation \auc{} between same- and cross-author text pairs, on the
same cohort, at token floors $1000$ and $2000$. \auc{} is the probability that a
same-author pair is scored more similar than a cross-author pair; $0.5$ is chance.
Table~\ref{tab:bakeoff} collects the results, with the prior study's count-based
baselines for reference---FastText \citep{bojanowski2017enriching},
\textsc{word2vec} \citep{mikolov2013distributed}, and Burrows's Delta
\citep{burrows2002delta}.

The count-based representations remain the strongest here: \textsc{word2vec}
($0.946$), FastText ($0.915$), and Delta ($0.907$) separate authors better than any
neural representation we test. The question this section answers is therefore not
whether a neural model wins document-level authorship---on $2.18$M tokens and about
eleven authors it does not---but which transfer route is the right inductive bias
for a zero-resource abjad, and why. The answer is a tokenizer-coverage effect the
numbers below make concrete.

\paragraph{A tokenizer-free byte encoder transfers.}
Off the shelf, with no Syriac training, CANINE encodes each form by mean-pooling
its codepoint encoding and reaches \auc{} $0.870$ / $0.849$. This is the first
demonstration that a pretrained multilingual byte model transfers to Syriac
authorship, and it already exceeds from-scratch byte- and character-level models
trained on the corpus.

\paragraph{Continued pretraining helps the LM, not authorship.}
Adapting CANINE to Syriac with LoRA ($1.06$M of $133$M parameters) lowers held-out
masked pseudo-bits-per-byte from $2.000$ to $1.859$ and raises masked-codepoint
accuracy from $0.307$ to $0.342$ relative to a frozen linear probe
(Table~\ref{tab:intrinsic}), so the encoder demonstrably learns Syriac. It does
\emph{not} improve authorship separation (\auc{} $0.857$ / $0.838$, slightly below
off the shelf): the masked objective is local, and document-mean pooling washes out
what it sharpens. The transfer win is the off-the-shelf encoder; the adaptation win
is the intrinsic objective.

\paragraph{Script coverage beats language count.}
Glot500, a subword model covering $500$+ languages including a nominal Syriac, does
worse than CANINE (\auc{} $0.798$ / $0.781$). The cause is concrete: its tokenizer
breaks Syriac into $7.05$ pieces per word with no Syriac subwords---character
fallback---whereas a model that actually covers the script should produce real
subwords. Nominal language coverage is not script coverage.

\paragraph{Cross-script transfer through Hebrew.}
Mapping Syriac into Hebrew script (a near $1{:}1$ abjad correspondence) and
encoding it with a Hebrew model (AlephBERT) reaches \auc{} $0.888$ / $0.857$, above
CANINE. The tokenizer confirms the mechanism: AlephBERT segments transliterated
Syriac into $2.51$ real subwords per word ($96.5\%$ covered), against Glot500's
$7.05$. This gives a clean ordering for a zero-resource abjad: char-fallback
multilingual ($0.798$) $<$ tokenizer-free byte ($0.870$) $<$ shared-script Semitic
transfer ($0.888$).

\paragraph{A supervised head matches, not beats, the count-based metric.}
Adding the prior study's leave-one-author-out supervised-contrastive head over the
encoder vectors gives the strongest \emph{neural} separation: CANINE $+$ head
$0.961\pm0.030$ at floor $1000$, Hebrew $+$ head $0.946\pm0.040$. This is on par
with the same head over the count-based FastText vectors ($0.966$ in the prior
study) rather than above it: the head, not the encoder, supplies the gain. Two
cautions attach to these numbers. The head trains on an $11$-author cohort and is nondeterministic across
runs even at a fixed seed, so we report a five-seed mean $\pm$ SD rather than a
single value; an earlier single run read $0.99$, which the sweep shows was an
optimistic draw. And at floor $2000$ the spread is large ($\pm0.10$ for CANINE), so
we do not read a single best pipeline there: CANINE and Hebrew with the head
overlap within their seed variation. The head is a real and general improvement
over raw cosine, but its margin on this cohort should be read as a range.

\begin{table}[t]
  \centering
  % Generated by neural/results.py -- do not edit by hand.
\begin{tabular}{lcc}
\toprule
Representation & AUC (floor 1000) & AUC (floor 2000) \\
\midrule
FastText (parent) & 0.915 & 0.900 \\
word2vec (parent) & 0.946 & 0.938 \\
Burrows's Delta (parent) & 0.907 & 0.940 \\
Glot500-m (char fallback) & 0.798 & 0.781 \\
CANINE-c off-the-shelf & 0.870 & 0.849 \\
\quad + AV head (LOAO) & 0.961$\pm$0.030 & 0.917$\pm$0.099 \\
Hebrew transliteration & 0.888 & 0.857 \\
\quad + AV head (LOAO) & 0.946$\pm$0.040 & 0.924$\pm$0.046 \\
\bottomrule
\end{tabular}

  \caption{Authorship separation \auc{} at token floors $1000$ / $2000$, same
  cohort and metric as the prior count-based study (its FastText, word2vec, and
  Delta numbers shown for reference). AV-head rows are five-seed mean $\pm$ SD; all
  other neural rows are deterministic given the encoder.}
  \label{tab:bakeoff}
\end{table}

\begin{table}[t]
  \centering
  % Generated by neural/results.py -- do not edit by hand.
\begin{tabular}{lrcc}
\toprule
 & trainable params & masked acc. & pseudo-bpb \\
\midrule
Frozen linear probe & 23,839 & 0.307 & 2.000 \\
LoRA continued-pretrain & 1,056,031 & 0.342 & 1.859 \\
\bottomrule
\end{tabular}

  \caption{Intrinsic language-model effect of LoRA continued pretraining against a
  frozen linear probe, isolated with an identical masked pseudo-bits-per-byte
  scorer (bidirectional, so comparable only among CANINE variants).}
  \label{tab:intrinsic}
\end{table}

% ===================================================================== %
\section{Factored Morphology and Textual Restoration}
\label{sec:other}

\paragraph{Explicit root/pattern factorisation does not help here.}
SEDRA's encoding splits every form into a consonant (root) channel and a vowel
(pattern) channel for free, which invites a two-stream encoder with explicit root
and pattern tiers. We compare it against a flat encoder reading the raw vocalised
characters, both trained with the same supervised-contrastive objective on the
root and evaluated by root-nearest-neighbour retrieval (Table~\ref{tab:factored}).
The factored model does not beat the flat one ($0.972$ vs.\ $0.978$ on seen roots,
$0.994$ vs.\ $0.994$ on unseen roots). Retrieval is near ceiling for both because
the consonantal skeleton is directly visible in the input, so surface overlap
already encodes the root; making the split explicit adds nothing when the split is
free. The factorisation would matter only where the channels are not separable by
inspection.

\paragraph{Character-level restoration.}
As an application demonstrator we train a causal character Transformer
($620$K parameters) on the DSC and evaluate lacuna restoration by synthetic
masking of held-out text. It restores masked characters at $0.443$ and whole spans
exactly at $0.089$, against a unigram floor near $0.19$ on the character metric,
and its fills are morphologically valid Syriac. A bidirectional masked objective at
this model and data scale collapses to the unigram prior; the causal objective,
which supervises every position, is what learns. This is a small but working Syriac
restorer in the spirit of neural restoration for other ancient languages
\citep{assael2022ithaca}.

\begin{table}[t]
  \centering
  % Generated by neural/results.py -- do not edit by hand.
\begin{tabular}{lrcc}
\toprule
 & params & seen-root NN & unseen-root NN \\
\midrule
Flat (vocalised chars) & 217,216 & 0.978 & 0.994 \\
Factored (root/pattern) & 251,536 & 0.972 & 0.994 \\
\midrule
$\Delta$ (factored $-$ flat) & & -0.007 & 0.000 \\
\bottomrule
\end{tabular}

  \caption{Factored root/pattern encoder vs.\ a flat encoder, root-nearest-neighbour
  retrieval on SEDRA forms split by root. ``Unseen'' roots are absent from training,
  the decisive generalisation case; both models are near ceiling.}
  \label{tab:factored}
\end{table}

% ===================================================================== %
\section{Discussion}
\label{sec:discussion}

The results divide along the line between intrinsic and document-level objectives,
the same division the prior count-based study found. Intrinsically, structure and
transfer both pay off: the vocaliser learns the pattern morpheme well enough to
generalise across register and to unseen skeletons, and continued pretraining
measurably improves the language-model objective. At the document level, the
picture is different. Continued pretraining does not help authorship, because
mean-pooling discards the local signal it sharpens; the supervised head helps most,
because it learns a metric directly on the pooled vectors. The practical
recommendation for a zero-resource abjad follows from the transfer ordering: prefer
a tokenizer-free encoder or transliteration into a related-language model's script
over a large multilingual model whose tokenizer misses the script, and add a
supervised head where author labels exist.

Two of our results are negative under their stated conditions, and both are
informative. Factoring root from pattern does not help when the skeleton is already
visible, which says the intrinsic root signal is easy and locates the difficulty
elsewhere. The supervised head's instability on an $11$-author cohort is a caution
about single-run headline numbers on small pools generally, not only here. The
broader method---inject structure through a morphological objective, transfer
through script rather than language count, and learn a metric for the document
task---should carry to other templatic, low-resource languages with the same
root-and-pattern, hapax-heavy profile.

% ===================================================================== %
\section{Limitations}
\label{sec:limitations}

The vocaliser's supervision is the New-Testament-scoped SEDRA lexicon; we measure
cross-register transfer to classical text but still train on one register. The
classical gold itself has limits we make explicit: scribes point partially, so we
score only pointed slots; one East-Syriac mark (U+073C) is genuinely ambiguous
between a vowel and the \emph{rukkaka} dot; and a few archaic variant letters pass
through the CAL map unconverted. The authorship cohort is small---about $11$
authors at the higher token floor, from a single corpus---so absolute \auc{} values
are cohort-specific and support relative comparison between representations rather
than population estimates; the supervised head's seed sensitivity is a direct
consequence of that scale, and its floor-$2000$ numbers are too noisy to rank. The
neural models are intentionally small and run on a single corpus; they probe
method under data scarcity, not an upper bound. Continued pretraining was tuned
lightly and not pushed; a different adaptation budget might change the intrinsic
numbers, though not the qualitative split between intrinsic and document-level
effects.

% ===================================================================== %
\section{Conclusion}
\label{sec:conclusion}

We treated the problem of building neural representations for a language with too
little text to pretrain on by injecting information from three other sources:
structure, through a vocaliser that restores Syriac pointing as the pattern
morpheme of a root-and-pattern system; transfer, through tokenizer-free and
cross-script encoders; and the corpus's own latent vocalisation signal, which we
turned into a classical gold set by deriving its mark-to-vowel mapping from data.
The vocaliser generalises across register and to unseen skeletons; cross-script
transfer through Hebrew outperforms both a byte encoder and a larger multilingual
model whose tokenizer misses the script; and a supervised head gives the best
authorship separation, within a seed range we report rather than a single number.
The negative results---no gain from explicit factorisation, no authorship gain from
continued pretraining---sharpen rather than weaken the account of where each kind of
information helps. The code, evaluation harness, and gold-set construction are
released, with SEDRA-derived data regenerated locally under its licence.

% ===================================================================== %
\section*{Reproducibility}
All numbers are produced by a single driver (\texttt{neural/results.py}) that runs
each experiment, writes a JSON manifest, and emits the table files used here, so no
value is transcribed by hand; the off-the-shelf and intrinsic numbers are
deterministic given the encoder, and the seed-sensitive AV-head numbers are
reported as five-seed mean $\pm$ SD. Experiments run on a single Apple-silicon
(MPS) device. The DSC \citep{digitalsyriaccorpus} and ETCBC corpora
\citep{etcbcsyriac} are openly licensed; the SEDRA lexicon \citep{kiraz1994sedra}
is licence-restricted, so we release the code that regenerates the derived data
from a user-provided SEDRA source rather than the data, and cite Kiraz as required.

% ===================================================================== %
\bibliographystyle{plainnat}
\bibliography{references}

\end{document}
